%% Question bank chapter file - generated by bank_tools.py build (do not edit while the bank is still being built:
%% it is regenerated from the per-test data in Question Bank\data). Every question is one bankq block:
%%   \begin{bankq}[id=..., source={...}, paper=.., qno=.., type=..., subtopic={...}, answer={...}, answer_src=key|solved,
%%                 working={...}, difficulty=1-10, diff_parts={K S I C F}, diff_note={...},
%%                 group=<passage id or empty>, also_in={duplicate copies in other tests}]
%% difficulty = min(10, K+S+I+C+F-1): K concepts 1-3, S steps 1-4, I insight 0-2, C computation 0-2, F format 0-1
%%   ...question LaTeX, options, figures (\bfig / \bfigi -> img/)...
%%   \end{bankq}
%% Passages shared by several questions are bankpassage blocks; their questions carry group=<passage id>.
\documentclass[10pt]{article}
\usepackage[a4paper,margin=1.8cm]{geometry}
\usepackage[T1]{fontenc}
\usepackage{amsmath,amssymb,mathtools}
\usepackage{graphicx}
\usepackage[export]{adjustbox}
\usepackage{xcolor,needspace,keyval}
\usepackage{enumitem}
\usepackage{booktabs,array,tabularx}
\usepackage[hidelinks]{hyperref}
\setlength{\parindent}{0pt}
\setlength{\parskip}{3pt}
\makeatletter
\define@key{bq}{id}{\def\bq@id{#1}}\define@key{bq}{source}{\def\bq@src{#1}}\define@key{bq}{paper}{\def\bq@paper{#1}}
\define@key{bq}{qno}{\def\bq@qno{#1}}\define@key{bq}{type}{\def\bq@type{#1}}\define@key{bq}{subtopic}{\def\bq@sub{#1}}
\define@key{bq}{answer}{\def\bq@ans{#1}}\define@key{bq}{answer_src}{\def\bq@asrc{#1}}\define@key{bq}{working}{\def\bq@work{#1}}
\define@key{bq}{group}{\def\bq@group{#1}}\define@key{bq}{also_in}{\def\bq@also{#1}}
\define@key{bq}{difficulty}{\def\bq@diff{#1}}\define@key{bq}{diff_parts}{\def\bq@dparts{#1}}\define@key{bq}{diff_note}{\def\bq@dnote{#1}}
\newenvironment{bankq}[1][]{\def\bq@work{}\def\bq@also{}\def\bq@group{}\def\bq@diff{?}\def\bq@dparts{}\def\bq@dnote{}\setkeys{bq}{#1}%
  \par\needspace{4\baselineskip}\medskip\noindent{\small\textbf{\bq@id}\hfill\color{gray}[\bq@src, Paper \bq@paper, Q\bq@qno]
  \bq@type\ $\cdot$ \bq@sub\ $\cdot$ \textbf{D\bq@diff}}\par\nopagebreak\relax}
  {\par{\small\color{teal}Answer: \textbf{\bq@ans}\ifx\bq@work\@empty\else\ (solved: \bq@work)\fi
  \ifx\bq@also\@empty\else\par\color{gray}Also in: \bq@also\fi}\par}
\define@key{bp}{id}{\def\bp@id{#1}}\define@key{bp}{source}{\def\bp@src{#1}}\define@key{bp}{questions}{\def\bp@qs{#1}}
\newenvironment{bankpassage}[1][]{\setkeys{bp}{#1}\par\needspace{4\baselineskip}\bigskip\noindent
  {\small\textbf{Passage \bp@id}\hfill\color{gray}[\bp@src; questions \bp@qs]}\par\nopagebreak\relax\itshape}{\par}
\makeatother
\newcommand{\bfig}[2]{\par\smallskip{\centering\includegraphics[width=#1pt,max width=0.62\linewidth,max height=0.32\textheight]{img/#2}\par}\smallskip}
\newcommand{\bfigi}[2]{\includegraphics[width=#1pt,max width=\linewidth,max height=0.32\textheight,valign=c]{img/#2}}
\newcommand{\bmissing}[1]{\textit{[figure missing in the source conversion: \detokenize{#1}]}}
\newenvironment{qoptions}{\begin{itemize}[leftmargin=1.8em,itemsep=1pt,topsep=2pt,label={}]}{\end{itemize}}
\newcommand{\opt}[2]{\item[(#1)] #2}
\providecommand{\fwstart}{}\providecommand{\fwend}{}\providecommand{\qfullsplit}{}
\newcolumntype{L}[1]{>{\raggedright\arraybackslash}p{#1}}
\newcolumntype{C}[1]{>{\centering\arraybackslash}p{#1}}
\newcolumntype{R}[1]{>{\raggedleft\arraybackslash}p{#1}}
\providecommand{\sectionhdr}[1]{}\providecommand{\subsectionhdr}[1]{}


\begin{document}
{\Large\bfseries Thermodynamics --- Question Bank}\par\smallskip
{\small Physics chapter P15. Questions from all converted test papers, grouped by source series. Answers: printed key where it matches the paper, otherwise solved (answer\_src=solved, with working). Exact repeats across papers appear once, with also\_in listing the other copies.}\par\medskip\hrule

\section*{Score-I papers}
\begin{bankq}[id=SI12-P1-PHY-S1-Q10, source={Score-I - Test 12 (09 Mar 2022)}, paper=1, qno=10, type=multi, subtopic={Processes}, answer={A,B}, answer_src=key, difficulty=4, diff_parts={K1 S2 I0 C1 F1}, diff_note={concepts: work in gas processes | steps: states A, B, C, D from the P-T graph; work in AB; work in BC with the ln term; work in CD and DA and the net work}]
Two moles of helium gas are taken over the cycle $ABCDA$, as shown in the $P$--$T$ diagram. (Use: $\ln2=0.7$)
\bfig{97.5}{SI12_p1_phy_s1_q10_1.png}
\begin{qoptions}
\opt{A}{The work done by the gas from $A$ to $B$ is $600R$}
\opt{B}{The work done by the gas from $B$ to $C$ is $840R$}
\opt{C}{The net work done by the gas in the cycle $ABCDA$ is $276R$}
\opt{D}{The net work done on the gas in the cycle $ABCDA$ is $1076R$}
\end{qoptions}
\end{bankq}

\section*{Score Advanced papers}
\begin{bankq}[id=SA01-P1-PHY-S2-Q08, source={Score Advanced - Test 01 (22 Sep 2020)}, paper=1, qno=8, type=multi, subtopic={Processes}, answer={A,B,C,D}, answer_src=key, difficulty=8, diff_parts={K2 S2 I2 C2 F1}, diff_note={concepts: insulated chamber with piston, processes | steps: equilibrium; adiabatic processes; final state; check}]
A thermally insulated chamber of volume 2V$_0$ is divided by a frictionless insulating piston of area S into two equal parts A and B. Part A has an ideal gas at pressure P$_0$ and temperature T$_0$ and in part B is vacuum. A massless spring of force constant K is connected with the piston and the wall of the container as shown. Initially the spring is unstretched. The ideal gas in chamber A is allowed to expand. Let in equilibrium the spring be compressed by a length x$_0$, then\par
\bfig{146.8}{SA01_p1_phy_s2_q08_1.png}
\begin{qoptions}
\opt{A}{Final pressure of the gas is Kx$_0$/S}
\opt{B}{Work done by the gas is $\frac12$ kx$_0^2$}
\opt{C}{Magnitude of change in internal energy of the gas is $\dfrac12\mathrm{kx}_0^2$}
\opt{D}{Temperature of the gas in equilibrium is less than initial temperature}
\end{qoptions}
\end{bankq}
\begin{bankq}[id=SA02-P1-PHY-S2-Q03, source={Score Advanced - Test 02 (13 May 2021)}, paper=1, qno=3, type=single, subtopic={Processes}, answer={C}, answer_src=solved, working={A to B isochoric heating, B to C isobaric expansion, C to A compression with heat loss; the cycle is clockwise}, difficulty=5, diff_parts={K2 S2 I1 C1 F0}, diff_note={concepts: thermodynamic signs on a p V diagram, processes | steps: identify each process; signs of work heat and energy; match}]
An ideal gas is taken along the reversible processes as represented by the adjoining diagram.\par
\bfig{99.9}{SA02_p1_phy_s2_q03_1.png}
\noindent{\footnotesize\setlength{\tabcolsep}{2pt}\begin{tabular}{|L{.7cm}|L{3.4cm}|L{.7cm}|L{2.2cm}|}\hline
 & List-I & & List-II\\\hline
(P) & For process B $\to$ C & (1) & $\Delta$Q $>$ 0\\\hline
(Q) & For process A $\to$ B & (2) & $\Delta$W $>$ 0\\\hline
(R) & For cycle A $\to$ B $\to$ C $\to$ A & (3) & $\Delta$U $>$ 0\\\hline
(S) & For process C $\to$ A & (4) & $\Delta$W = 0\\\hline
 & & (5) & $\Delta$Q $<$ 0\\\hline
\end{tabular}}\par
\begin{qoptions}
\opt{A}{P $\to$ 1,2,3;Q $\to$ 1,2,5;R $\to$ 1,2,3;S $\to$ 1,2}
\opt{B}{P $\to$ 1,2;Q $\to$ 1,3,4;R $\to$ 1,4,5;S $\to$ 5}
\opt{C}{P $\to$ 1,2,3;Q $\to$ 1,3,4;R $\to$ 1,2;S $\to$ 5}
\opt{D}{P $\to$ 1,3;Q $\to$ 1,3,4;R $\to$ 5;S $\to$ 4,5}
\end{qoptions}
\end{bankq}
\begin{bankq}[id=SA02-P2-PHY-S1-Q05, source={Score Advanced - Test 02 (13 May 2021)}, paper=2, qno=5, type=multi, subtopic={Processes}, answer={A,B,C}, answer_src=solved, working={TV\^{}(gamma-1)=constant gives m=V0/(T0(gamma-1)); CV=R/(gamma-1); CP=CV+R}, difficulty=7, diff_parts={K2 S2 I2 C1 F1}, diff_note={concepts: slope of an adiabatic V-T curve, processes | steps: adiabatic relation; differentiate; specific heats}]
A gas is undergoing an adiabatic process. At a certain stage, the values of volume and temperature are V$_0$ and T$_0$ respectively and the modulus of slope of V--T curve is m. Choose the correct options-
\begin{qoptions}
\opt{A}{The value of C$_V$ = $\dfrac{RT_0 m}{V_0}$}
\opt{B}{The value of C$_P$ = $\left(1 + \dfrac{T_0 m}{V_0}\right)$R}
\opt{C}{m = $\dfrac{V_0}{T_0(\gamma - 1)}$}
\opt{D}{m = $\dfrac{-V_0}{T_0(\gamma - 1)}$}
\end{qoptions}
\end{bankq}
\begin{bankq}[id=SA05-P1-PHY-S1-Q07, source={Score Advanced - Test 05 (15 Jul 2021)}, paper=1, qno=7, type=single, subtopic={Heat engines}, answer={A}, answer_src=key, difficulty=3, diff_parts={K1 S1 I1 C1 F0}, diff_note={concepts: direction of a cycle on a P-V diagram | steps: clockwise cycle is an engine; anticlockwise is a refrigerator}]
If the P-V diagrams of two thermodynamics devices working in a cyclic process are as shown in the figure, then :-\par
\bfig{122.7}{SA05_p1_phy_s1_q07_1.png}
\begin{qoptions}
\opt{A}{A is a heat engine, B is a heat pump/refrigerator}
\opt{B}{B is a heat engine, A is a heat pump/refrigerator}
\opt{C}{both A and B are heat engines}
\opt{D}{both A and B are heat pumps/refrigerator}
\end{qoptions}
\end{bankq}
\begin{bankq}[id=SA07-P2-PHY-S1-Q01, source={Score Advanced - Test 07 (16 Apr 2023)}, paper=2, qno=1, type=multi, subtopic={Processes}, answer={A,B,C,D}, answer_src=key, difficulty=9, diff_parts={K3 S2 I2 C2 F1}, diff_note={concepts: cyclic process of a helium gas, processes, adiabatic and isothermal relations | steps: adiabatic and isothermal relations; work in each step; efficiency}]
1 mol He gas undergoes a cyclic process ABCA as shown in the figure. Process AB, BC \& CA are isothermal, isochoric and adiabatic respectively. Temperature of gas at point A is T$_0$. On the basis of given information identify correct statements. (2$^{2/3}$ = 1.5) and (ln2 = 0.7)\par
\bfig{97.3}{SA07_p2_phy_s1_q01_1.png}
\begin{qoptions}
\opt{A}{Change in internal energy in process BC is --RT$_0$/2.}
\opt{B}{Work done by gas in process CA is --RT$_0$/2.}
\opt{C}{Work done by gas in cyclic process is RT$_0$/5.}
\opt{D}{Efficiency of cyclic process is 2/7.}
\end{qoptions}
\end{bankq}
\begin{bankq}[id=SA09-P1-PHY-S2-Q09, source={Score Advanced - Test 09 (29 Apr 2023)}, paper=1, qno=9, type=multi, subtopic={First law}, answer={A,B,C}, answer_src=key, difficulty=8, diff_parts={K3 S1 I2 C2 F1}, diff_note={concepts: work and temperature of a gas between charged pistons, first law, electrostatic force between pistons | steps: electrostatic force between pistons; first law}]
An ideal gas having degree of freedom f is kept inside a thermally insulated vessel. Vessel having two large pistons each with charge of +q and $-$q. Assume charges are uniformly distributed. At an instant charge at pistons is increased by $\beta$ times with help of an external agent. Then which of the following statements are correct ? No friction between wall and piston.\par
\bfig{108.3}{SA09_p1_phy_s2_q09_1.png}
\begin{qoptions}
\opt{A}{Work done by the agent in the process is ($\beta^2-1$)nRT$_1$ where T$_1$ is initial temp.}
\opt{B}{Work done by the agent in the process is ($\beta^2-1$)$\dfrac{\text{q}^2\text{L}_1}{\text{A}\in_0}$ where L$_1$ is initial separation between pistons}
\opt{C}{The ratio of final to initial temperature of gas in equilibrium is $\dfrac{2\beta^2+\text{f}}{2+\text{f}}$}
\opt{D}{The ratio of final to initial temperature of gas in equilibrium is $\dfrac{\beta^2+\text{f}}{2+\text{f}}$}
\end{qoptions}
\end{bankq}
\begin{bankq}[id=SA09-P1-PHY-S3-Q04, source={Score Advanced - Test 09 (29 Apr 2023)}, paper=1, qno=4, type=numerical, subtopic={Heat engines}, answer={2.5}, answer_src=key, difficulty=8, diff_parts={K3 S1 I2 C2 F1}, diff_note={concepts: temperature rise of the liquid per cycle, heat engines, net work per cycle from areas | steps: net work per cycle from areas; heat lost by the gas equals the heat gained by the liquid}]
In the given figure two co-axial cylinders are placed as shown. Inner cylinder of radius r contains one mole of an ideal gas enclosed by a piston which can move in inner cylinder without friction and the space between the cylinder of radius 2r and radius r contains liquid. Heat is supplied to the gas kept in inner cylinder at a constant rate of 100 R J/cycle by some external agent. The cyclic process for the gas is represented by the indicator diagram as shown. Now if heat lost by the gas equal to heat gain by the liquid and radiation is negligible, find the rate of increasing of temperature of liquid per cycle (mass of liquid is 1 kg and specific heat capacity = 10 R in SI unit, temperature at points 1, 2 and 3 in cycle are 300 K, 400 K and 600 K respectively, R is universal gas constant).\par
\bfig{245.2}{SA09_p1_phy_s3_q04_1.png}
\end{bankq}
\begin{bankq}[id=SA09-P2-PHY-S3-Q04, source={Score Advanced - Test 09 (29 Apr 2023)}, paper=2, qno=4, type=numerical, subtopic={Processes}, answer={50}, answer_src=key, difficulty=8, diff_parts={K3 S1 I2 C2 F1}, diff_note={concepts: volume where a process changes from endothermic to exothermic, processes, dQ=0 on the straight line process | steps: dQ=0 on the straight line process}]
One mole of an ideal mono-atomic gas undergoes a process A $\to$ B that is a straight line on PV indicator diagram of the process as shown in the figure. Find volume of the gas in litre when the process turns from an endothermic to an exothermic one.\par
\bfig{145.0}{SA09_p2_phy_s3_q04_1.png}
\end{bankq}
\begin{bankq}[id=SA10-P1-PHY-S2-Q08, source={Score Advanced - Test 10 (06 May 2023)}, paper=1, qno=8, type=single, subtopic={Processes}, answer={D}, answer_src=key, difficulty=7, diff_parts={K3 S1 I2 C2 F0}, diff_note={concepts: work, internal energy changes in four processes, areas under the curves and temperatures relative to the isotherm | steps: areas under the curves and temperatures relative to the isotherm}]
Volume versus pressure curves for n moles of an ideal gas are given for four process as shown in figure. AB is an isothermal curve. Match the entries of List-I with the enteries of List-II.\par
\bfig{108.3}{SA10_p1_phy_s2_q08_1.png}
{\footnotesize\setlength{\tabcolsep}{2pt}\noindent\begin{tabular}{|l|l|l|L{4.2cm}|}
\hline
 & List-I & & \multicolumn{1}{c|}{List-II}\\ \hline
(I) & For process 1 & (P) & Work done is zero.\\ \hline
(II) & For process 2 & (Q) & Work done is maximum.\\ \hline
(III) & For process 3 & (R) & Internal energy will increase.\\ \hline
(IV) & For process 4 & (S) & Internal energy will decrease.\\ \hline
 & & (T) & Internal energy remains constant.\\ \hline
\end{tabular}}\par
\begin{qoptions}
\opt{A}{I $\to$ P,R;II $\to$ R;III $\to$ S;IV $\to$ P,S}
\opt{B}{I $\to$ Q,R;II $\to$ R;III $\to$ S;IV $\to$ P,Q}
\opt{C}{I $\to$ Q,R;II $\to$ S;III $\to$ R;IV $\to$ P,S}
\opt{D}{I $\to$ Q,R;II $\to$ R;III $\to$ S;IV $\to$ P,S}
\end{qoptions}
\end{bankq}
\begin{bankq}[id=SA11-P1-PHY-S1-Q06, source={Score Advanced - Test 11 (09 May 2023)}, paper=1, qno=6, type=multi, subtopic={Processes}, answer={A,B,C,D}, answer_src=key, difficulty=8, diff_parts={K3 S1 I2 C2 F1}, diff_note={concepts: monatomic ideal gas in a given process, processes, work and heat from the process equation | steps: work and heat from the process equation}]
One mole of an ideal monatomic gas has initial temperature T$_0$, is made to go through the cycle abca shown in the given figure. If U denotes the internal energy, then choose the correct alternative.\par
\bfig{105.6}{SA11_p1_phy_s1_q06_1.png}
\begin{qoptions}
\opt{A}{U$_\text{c}$ $>$ U$_\text{b}$ $>$ U$_\text{a}$}
\opt{B}{U$_\text{c}$ $-$ U$_\text{b}$ = 3RT$_0$}
\opt{C}{U$_\text{c}$ $-$ U$_\text{a}$ = $\dfrac{9\text{RT}_0}{2}$}
\opt{D}{U$_\text{b}$ $-$ U$_\text{a}$ = $\dfrac{3\text{RT}_0}{2}$}
\end{qoptions}
\end{bankq}
\begin{bankpassage}[id=SA13-P2-PHY-S2-G9-10, source={Score Advanced - Test 13 (13 May 2023)}, questions={9-10}]
Two vertical cylinders of different areas of cross-sections (A$_2$ $<$ A$_1$) shown in the figures have pistons of same mass (m$_1$ = m$_2$) fitted on the top. The cylinders enclose N moles of monatomic gas. The initial pressure and initial volume of the gas are P and V, respectively. It is given that spring I is relaxed. Now through the electric heater, heat is slowly supplied. Finally spring II becomes relaxed. It is given that the cylinders and the pistons are thermally insulated.\par
\bfig{160.9}{SA13_p2_phy_s2_para9_10_1.png}
\end{bankpassage}
\begin{bankq}[id=SA13-P2-PHY-S2-Q09, source={Score Advanced - Test 13 (13 May 2023)}, paper=2, qno=9, type=single, subtopic={First law}, answer={B}, answer_src=key, difficulty=7, diff_parts={K3 S1 I2 C2 F0}, diff_note={concepts: final pressure of gas in two cylinders, first law, energy balance for adiabatic processes | steps: energy balance for adiabatic processes}, group=SA13-P2-PHY-S2-G9-10]
What is the final pressure?
\begin{qoptions}
\opt{A}{$\text{P}+\text{mg}\left(\dfrac{\text{A}_1+\text{A}_2}{\text{A}_1\text{A}_2}\right)$}
\opt{B}{$\text{P}+\text{mg}\left(\dfrac{\text{A}_1-\text{A}_2}{\text{A}_1\text{A}_2}\right)$}
\opt{C}{$\text{P}+2\text{mg}\left(\dfrac{\text{A}_1-\text{A}_2}{\text{A}_1\text{A}_2}\right)$}
\opt{D}{$\text{P}-2\text{mg}\left(\dfrac{\text{A}_1+\text{A}_2}{\text{A}_1\text{A}_2}\right)$}
\end{qoptions}
\end{bankq}
\begin{bankq}[id=SA13-P2-PHY-S2-Q10, source={Score Advanced - Test 13 (13 May 2023)}, paper=2, qno=10, type=single, subtopic={First law}, answer={B}, answer_src=key, difficulty=7, diff_parts={K3 S1 I2 C2 F0}, diff_note={concepts: final volume of gas in the cylinder, processes, process equation | steps: process equation}, group=SA13-P2-PHY-S2-G9-10]
What is the final volume?
\begin{qoptions}
\opt{A}{$\text{V}+\dfrac{\text{mg}}{\text{k}}\dfrac{(\text{A}_1+\text{A}_2)(\text{A}_1^2+\text{A}_2^2)}{\text{A}_1\text{A}_2}$}
\opt{B}{$\text{V}+\dfrac{\text{mg}}{\text{k}}\dfrac{(\text{A}_1-\text{A}_2)(\text{A}_1^2+\text{A}_2^2)}{\text{A}_1\text{A}_2}$}
\opt{C}{$\text{V}+\dfrac{\text{mg}}{2\text{k}}\dfrac{(\text{A}_1-\text{A}_2)(\text{A}_1^2+\text{A}_2^2)}{\text{A}_1\text{A}_2}$}
\opt{D}{$\text{V}+\dfrac{2\text{mg}}{\text{k}}\dfrac{(\text{A}_1-\text{A}_2)(\text{A}_1^2+\text{A}_2^2)}{\text{A}_1\text{A}_2}$}
\end{qoptions}
\end{bankq}
\begin{bankpassage}[id=SA13-P2-PHY-S4-G5-6, source={Score Advanced - Test 13 (13 May 2023)}, questions={5-6}]
A Carnot heat engine draws heat from a reservoir at temperature T$_1$ and rejects heat to another reservoir at temperature T$_3$. The Carnot forward cycle engine drives a Carnot reversed cycle engine or Carnot refrigerator which absorbs heat from reservoir at temperature T$_2$ and rejects heat to a reservoir at temperature T$_3$. The high temperature T$_1$ is 600 K and low temperature T$_2$ is 300 K.\par
\end{bankpassage}
\begin{bankq}[id=SA13-P2-PHY-S4-Q05, source={Score Advanced - Test 13 (13 May 2023)}, paper=2, qno=5, type=numerical, subtopic={Heat engines}, answer={400}, answer_src=key, difficulty=8, diff_parts={K3 S1 I2 C2 F1}, diff_note={concepts: temperature of the middle reservoir of coupled Carnot machines, heat engines, equal heat flows | steps: equal heat flows; efficiency and COP}, group=SA13-P2-PHY-S4-G5-6]
If heat supplied to engine Q$_1$ is equal to heat absorbed by the refrigerator Q$_2$, the value of T$_3$ (in K) is :-\par
\end{bankq}
\begin{bankq}[id=SA13-P2-PHY-S4-Q06, source={Score Advanced - Test 13 (13 May 2023)}, paper=2, qno=6, type=numerical, subtopic={Heat engines}, answer={3}, answer_src=key, difficulty=8, diff_parts={K3 S1 I2 C2 F1}, diff_note={concepts: coefficient of performance of the Carnot refrigerator, heat pumps, T/(T hot - T cold) | steps: T/(T hot - T cold)}, group=SA13-P2-PHY-S4-G5-6]
The COP of carnot refrigerator is :-\par
\end{bankq}
\begin{bankq}[id=SA14-P1-PHY-S3-Q01, source={Score Advanced - Test 14 (16 May 2023)}, paper=1, qno=1, type=numerical, subtopic={Heat engines}, answer={4}, answer_src=key, difficulty=8, diff_parts={K3 S1 I2 C2 F1}, diff_note={concepts: cyclic process on an ideal gas, heat engines, work from the areas | steps: work from the areas; relate pressures and volumes}, also_in={Score Advanced - Test 20 (06 Apr 2024) P1 Q1}]
Figure shows a cyclic process ABCDBEA performed on an ideal gas. If P$_\text{A}$ = 2 atm.,P$_\text{B}$ =4.5 atm. and P$_\text{C}$ = 6 atm.; V$_\text{E}$ - V$_\text{A}$ = 0.05 litre, find the magnitude of work done by the gas (in J) in the complete process in nearest integer.\par
(1 atm, pressure =1x10$^5$ pa)\par
\bfig{112.4}{SA14_p1_phy_s3_q01_1.png}
\end{bankq}
\begin{bankq}[id=SA15-P1-PHY-S2-Q07, source={Score Advanced - Test 15 (18 May 2023)}, paper=1, qno=7, type=multi, subtopic={Processes}, answer={B,D}, answer_src=key, difficulty=8, diff_parts={K3 S1 I2 C2 F1}, diff_note={concepts: adiabatic partition in a ring-shaped tube, processes, adiabatic compression of one chamber by the heated gas | steps: adiabatic compression of one chamber by the heated gas}]
A ring shaped tube of volume $2V_0$ contains an ideal gas with $\gamma=1.5$. It has an adiabatic fixed wall S$_1$ and a movable adiabatic partition S$_2$, and a heater H. S$_2$ moves a quarter of the circle.\par
\bfig{121.0}{SA15_p1_phy_s2_q07_1.png}
\begin{qoptions}
\opt{A}{Final pressure in chamber 2 is $2P_0$}
\opt{B}{Final pressure in chamber 2 is $2\sqrt2\,P_0$}
\opt{C}{Work done on the gas in chamber 2 by the gas in chamber 1 is $(2\sqrt2-1)P_0V_0$}
\opt{D}{Heat supplied is $4P_0V_0(2\sqrt2-1)$}
\end{qoptions}
\end{bankq}
\begin{bankq}[id=SA17-P1-PHY-S3-Q04, source={Score Advanced - Test 17 (25 May 2023)}, paper=1, qno=4, type=numerical, subtopic={Processes}, answer={0.50}, answer_src=key, difficulty=7, diff_parts={K3 S1 I2 C1 F1}, diff_note={concepts: heat supplied in a process with P proportional to V, processes, work and internal energy for P/V constant | steps: work and internal energy for P/V constant}]
One mole of an ideal monatomic gas at temperature $T_0$ expands slowly according to the law\par
$P/V=$ constant. If the final temperature is $2T_0$, heat supplied to the gas is $\dfrac{RT_0}{x_0}$. Then find value of x?
\end{bankq}
\begin{bankq}[id=SA18-P2-PHY-S3-Q04, source={Score Advanced - Test 18 (27 May 2023)}, paper=2, qno=4, type=numerical, subtopic={Processes}, answer={3}, answer_src=key, difficulty=8, diff_parts={K3 S1 I2 C2 F1}, diff_note={concepts: heat in two successive gas processes, constant volume then constant pressure, work and heat | steps: constant volume then constant pressure; work and heat}]
Two moles of a diatomic gas undergoes two process in succession. Starting from temperature $T_1$, it is first heated at constant volume until its pressure is tripled. It is further heated at constant pressure until its volume is doubled. What is the average molar heat capacity for the whole process? If your answer is $nR$, where $R$ is universal gas constant. Round off $n$ to nearest integer.
\end{bankq}
\begin{bankq}[id=SA19-P1-PHY-S2-Q08, source={Score Advanced - Test 19 (04 Apr 2024)}, paper=1, qno=8, type=multi, subtopic={Processes}, answer={A,B,D}, answer_src=key, difficulty=9, diff_parts={K3 S2 I2 C2 F1}, diff_note={concepts: straight line path in the pressure volume diagram, processes, work above the isotherm | steps: work above the isotherm; T-V relation; maximum temperature}]
An ideal gas is taken from the state A ($P$, $V$) to the state B($P/2$, $2V$) along a straight line path as shown in figure. Select the correct statement from the following\par
\bfig{118.2}{SA19_p1_phy_s2_q08_1.png}
\begin{qoptions}
\opt{A}{Work done by the gas in going from A to B exceeds the work done in going from A to B under isothermal conditions,}
\opt{B}{In the T-V diagram, part AB would becomes a parabola}
\opt{C}{In the P-T diagram, path AB would be part of hyperbola}
\opt{D}{In going from A to B, the temp. T of gas first increases to max value and then decreases.}
\end{qoptions}
\end{bankq}
\begin{bankq}[id=SA23-P1-PHY-S2-Q06, source={Score Advanced - Test 23 (14 Apr 2024)}, paper=1, qno=6, type=multi, subtopic={Processes}, answer={C,D}, answer_src=key, difficulty=8, diff_parts={K3 S1 I2 C2 F1}, diff_note={concepts: gas between two piston cylinder arrangements connected by a valve, processes, first law for adiabatic or isothermal mixing | steps: first law for adiabatic or isothermal mixing}]
We have two piston cylinder arrangement connected by a valve as shown in figure.\par
\bfig{208.0}{SA23_p1_phy_s2_q06_1.png}
Same ideal monoatomic gas is filled in both cylinders at same temperature $T_0$ but at different pressures $2P_0$ and $P_0$. Both cylinders are connected by a thermal tube fitted with a valve, which is initially closed. Now valve is opened and gas starts leaking from cylinder A to B. Mr. Swapnil keeps pushing cylinder A and maintain $2P_0$ pressure in cylinder A while Mr. Sri holds the piston B at its original position during the motion. Process is stopped when pressure in both cylinders becomes $2P_0$. (Neglect any loss of heat in surrounding)
\begin{qoptions}
\opt{A}{Final temperature of gas in both cylinders is $\dfrac{19}{15}T_0$.}
\opt{B}{Final volume of cylinder A is $\dfrac{3}{10}$V.}
\opt{C}{Final temperature of gas in both cylinders $\dfrac{17}{15}T_0$.}
\opt{D}{Final volume of cylinder A is $\dfrac{7}{10}V_0$.}
\end{qoptions}
\end{bankq}
\begin{bankq}[id=SA24-P1-PHY-S2-Q08, source={Score Advanced - Test 24 (17 Apr 2024)}, paper=1, qno=8, type=single, subtopic={Processes}, answer={A}, answer_src=key, difficulty=7, diff_parts={K3 S1 I2 C2 F0}, diff_note={concepts: heat and work for a process P=P0+alpha/V, processes, work by integration | steps: work by integration; molar heat capacity}]
One mole of an ideal gas whose adiabatic exponent equals $\gamma$ undergoes a process $P=P_0+\alpha/v$, where $P_0$ and $\alpha$ are positive constants. Its volume increases from $v_1$ to $v_2$.\par
{\footnotesize\setlength{\tabcolsep}{1.5pt}\noindent\begin{tabular}{|c|L{3.0cm}|c|L{3.2cm}|} \hline & \multicolumn{1}{c|}{List-I} & & \multicolumn{1}{c|}{List-II}\\ \hline
(I) & Heat capacity of gas as a function of its volume & (P) & $\dfrac{P_0(V_2-V_1)}{\gamma-1}$\\ \hline
(II) & Internal energy increment of a gas & (Q) & $\dfrac{P_0(V_2-V_1)}{\gamma-1}+\alpha\ln\dfrac{V_2}{V_1}$\\ \hline
(III) & The work performed by it & (R) & $P_0(V_2-V_1)+\alpha\ln\dfrac{V_2}{V_1}$\\ \hline
(IV) & The amount of heat transferred & (S) & $\dfrac{\gamma R}{\gamma-1}+\dfrac{\alpha R}{p_0v}$\\ \hline
\end{tabular}}\par
\begin{qoptions}
\opt{A}{I $\rightarrow$ S;II $\rightarrow$ P;III $\rightarrow$ R;IV $\rightarrow$ Q}
\opt{B}{I $\rightarrow$ P;II $\rightarrow$ Q;III $\rightarrow$ R;IV $\rightarrow$ S}
\opt{C}{I $\rightarrow$ P;II $\rightarrow$ S;III $\rightarrow$ R;IV $\rightarrow$ Q}
\opt{D}{I $\rightarrow$ S;II $\rightarrow$ P;III $\rightarrow$ Q;IV $\rightarrow$ S}
\end{qoptions}
\end{bankq}
\begin{bankq}[id=SA24-P2-PHY-S2-Q09, source={Score Advanced - Test 24 (17 Apr 2024)}, paper=2, qno=9, type=multi, subtopic={Processes}, answer={A,C}, answer_src=key, difficulty=8, diff_parts={K3 S1 I2 C2 F1}, diff_note={concepts: graph of a thermodynamic cycle in other variables, processes, identify each process in the new diagram | steps: identify each process in the new diagram}]
An ideal gas under goes a thermodynamic cycle as shown in figure. Which of the following graphs represents the same cycle ?\par
\bfig{77.4}{SA24_p2_phy_s2_q09_1.png}
\begin{qoptions}
\opt{A}{\bfigi{70.3}{SA24_p2_phy_s2_q09_A.png}}
\opt{B}{\bfigi{68.6}{SA24_p2_phy_s2_q09_B.png}}
\opt{C}{\bfigi{70.8}{SA24_p2_phy_s2_q09_C.png}}
\opt{D}{\bfigi{71.3}{SA24_p2_phy_s2_q09_D.png}}
\end{qoptions}
\end{bankq}
\begin{bankq}[id=SA24-P2-PHY-S3-Q01, source={Score Advanced - Test 24 (17 Apr 2024)}, paper=2, qno=1, type=numerical, subtopic={Processes}, answer={5}, answer_src=key, difficulty=8, diff_parts={K3 S1 I2 C2 F1}, diff_note={concepts: heating of gas under a movable piston, processes, isobaric process | steps: isobaric process; work and heat}]
A vertical hollow cylinder contains an ideal gas. The gas is enclosed by a 5 kg movable piston with area of cross-section $5\times10^{-3}$ m$^2$. The gas is slowly heated from 300 K to 350 K and the piston rises by 0.1 m. The piston is now clamped at this position and the gas is cooled back to 300 K. The difference between the heat energy added during heating process and heat energy lost during cooling process is 11x. Find value of x. [Take $p_0=10^5$ N/m$^2$]
\end{bankq}
\begin{bankq}[id=SA25-P1-PHY-S2-Q07, source={Score Advanced - Test 25 (21 Apr 2024)}, paper=1, qno=7, type=multi, subtopic={First law}, answer={A,B,C}, answer_src=key, difficulty=8, diff_parts={K3 S1 I2 C2 F1}, diff_note={concepts: air entering an insulated cylinder through a valve, first law, work done by outside air | steps: work done by outside air; adiabatic}, also_in={Score Advanced - Test 26 (21 Apr 2024) P1 Q7}]
At the end of the thermally insulated cylinder with the piston there is a valve (figure) that covers a small hole, which opens and begins to let air from outside into the cylinder at a pressure difference $\Delta P=P_0/3$ ($P_0$ is atmospheric pressure). The valve does not allow air to escape from the cylinder to the outside. At the initial moment of time, the piston is pressed against the end of the cylinder, there is no air inside.\par
In the first case, the cylinder is filled with air to volume $V_0$, slowly moving the piston, after which it is stopped, and then the piston is released.\par
In the second case, the piston is abruptly moved to a position in which the volume under the piston is equal to $V_0$, so that the air does not have time to penetrate through the valve into the cylinder. In this position, fix the piston, wait until the cylinder is filled with air, just as in the first case, release the piston. Out side the cylinder is air at atmospheric pressure $P_0$ and temperature $T_0$. The friction of the piston against the walls, the mass of the piston, as well as the heat exchange of air with the piston and the cylinder walls can be neglected. Air can be considered a diatomic ideal gas. After the piston is released, the valve remains closed all the time.\par
\bfig{56.0}{SA25_p1_phy_s2_q06_2.png}
\begin{qoptions}
\opt{A}{For first case the temperature of air inside the cylinder at the moment the piston is stopped at volume $V_0$ is $T_0$.}
\opt{B}{For first case, the temperature of air inside the cylinder when the piston stops movement is $\dfrac87T_0$.}
\opt{C}{For second case, the temperature of the air in the cylinder after stopping of piston and filling with air is $\dfrac75T_0$.}
\opt{D}{For second case, the temperature of the air after piston is released and it stops movement is $\dfrac95T_0$.}
\end{qoptions}
\end{bankq}
\begin{bankq}[id=SA25-P1-PHY-S3-Q11, source={Score Advanced - Test 25 (21 Apr 2024)}, paper=1, qno=11, type=single, subtopic={Processes}, answer={C}, answer_src=key, difficulty=7, diff_parts={K3 S1 I2 C2 F0}, diff_note={concepts: heat given in a process shown on a V-T graph, processes, P from the slope | steps: P from the slope; first law}, also_in={Score Advanced - Test 26 (21 Apr 2024) P1 Q11}]
V$-$T graph of a thermodynamic process is shown in figure. If heat given to process is taken to be positive :-\par
($\Delta$Q is heat given to the gas in the process, $\Delta$U is change in internal energy in the process, $\Delta$W is the work done by the gas in the process).\par
\bfig{133.3}{SA25_p1_phy_s3_q11_1.png}
{\footnotesize\setlength{\tabcolsep}{1.5pt}\noindent\begin{tabular}{|c|L{2.0cm}|c|L{2.0cm}|} \hline & \multicolumn{1}{c|}{List$-$I} & & \multicolumn{1}{c|}{List$-$II}\\ \hline
(P) & JK & (1) & $\Delta W\geq0$\\ \hline
(Q) & KL & (2) & $\Delta Q<0$\\ \hline
(R) & LM & (3) & $\Delta W<0$\\ \hline
(S) & MJ & (4) & $\Delta Q>0$\\ \hline
 & & (5) & $\Delta U=0$\\ \hline
\end{tabular}}\par
\begin{qoptions}
\opt{A}{P $\rightarrow$ 2,3,5;Q $\rightarrow$ 4;R $\rightarrow$ 1,4,5;S $\rightarrow$ 2,3}
\opt{B}{P $\rightarrow$ 1,3,5;Q $\rightarrow$ 4;R $\rightarrow$ 1,4,5;S $\rightarrow$ 2,5}
\opt{C}{P $\rightarrow$ 2,3,5;Q $\rightarrow$ 1,4;R $\rightarrow$ 1,4,5;S $\rightarrow$ 2,3}
\opt{D}{P $\rightarrow$ 1,4;Q $\rightarrow$ 2,4;R $\rightarrow$ 1,3,5;S $\rightarrow$ 2,3}
\end{qoptions}
\end{bankq}
\begin{bankq}[id=SA27-P1-PHY-S3-Q06, source={Score Advanced - Test 27 (24 Apr 2024)}, paper=1, qno=6, type=numerical, subtopic={Processes}, answer={6}, answer_src=key, difficulty=8, diff_parts={K3 S1 I2 C2 F1}, diff_note={concepts: monatomic gas process with a given molar heat capacity, processes, C=Cv+R/(1-x) gives the polytropic exponent | steps: C=Cv+R/(1-x) gives the polytropic exponent; work}]
One mole of a monatomic ideal gas follows a process A $\rightarrow$ B, as shown. The specific heat of the process is $\dfrac{13\mathrm R}{6}$. The value of `n' on P-axis (see following figure) is\par
\bfig{108.0}{SA27_p1_phy_s3_q06_1.png}
\end{bankq}
\begin{bankq}[id=SA27-P2-PHY-S3-Q04, source={Score Advanced - Test 27 (24 Apr 2024)}, paper=2, qno=4, type=numerical, subtopic={Heat engines}, answer={5}, answer_src=solved, working={n R = 0.4 x 50/6 = 10/3; V1 = nRT1/P1 = 1e-3 m\^{}3. State 2: P2 = 2e6, V2 = 1e-3. Adiabatic: V3 = V1 2\^{}(3/5) = 1.5e-3 m\^{}3 at P = 1e6. W23 = (2000-1500)/(gamma-1) = 750 J, W31 = 1e6 x (1e-3-1.5e-3) = -500 J. Net heat = net work = 250 J; N/50 = 5.}, difficulty=8, diff_parts={K3 S1 I2 C2 F1}, diff_note={concepts: net heat of a cycle with isochoric adiabatic and isobaric steps, heat engines, state volumes from nRT/P | steps: state volumes from nRT/P; work in each step}]
A heat engine takes 0.4 moles of an ideal gas around the cycle as shown in the figure. Process 1 $\rightarrow$ 2 is isochoric, process 2 $\rightarrow$ 3 is adiabatic and process 3 $\rightarrow$ 1 is isobaric. The value of $\gamma$ for the gas is 5/3. Find the net heat flow during the cycle in joule. If your answer is N fill value of N/50 ?\par
[Take: Universal gas constant (R) = 50 / 6 J mol$^{-1}$K$^{-1}$ : 1 atm = 10$^6$ N/m$^2$ and 2$^{3/5}$ = 1.5]\par
\bfig{114.8}{SA27_p2_phy_s3_q04_1.png}
\end{bankq}

\section*{TOAS papers}
\begin{bankq}[id=TOAS12-P1-PHY-S2-Q08, source={TOAS - Test 12 (10 Jul 2022)}, paper=1, qno=8, type=numerical, subtopic={Heat engines}, answer={0.48}, answer_src=solved, working={TB = 4TA gives VB = 4VA; the adiabat BC gives TC = 2TA; Qin = nCp (3TA) = 7.5 nRTA; heat out = nCv TA + nRTA ln(8 sqrt2) = (1.5 + 2.43) nRTA; efficiency 1 - 3.93/7.5 = 0.48}, difficulty=5, diff_parts={K1 S3 I1 C0 F1}, diff_note={concepts: efficiency of a cycle with isobaric and adiabatic and isochoric and isothermal steps | steps: temperature and volume ratios from the data; temperature at C from the adiabat; heat in along AB; heat out along CD; heat out along DA; efficiency | I: use the adiabat to find TC}]
Helium is used as working substance in an engine working on the cycle as shown in the figure. Processes AB, BC, CD and DA are isobaric, adiabatic, isochoric and isothermal respectively. The ratio of maximum to minimum volume of helium during the cycle is $8\sqrt2$ and that of maximum to minimum absolute temperature is 4. Calculate efficiency ($\eta$) of the cycle ($\eta<1$).
\bfig{88.2}{TOAS12_p1_phy_s2_q08_1.png}
\end{bankq}
\begin{bankq}[id=TOAS12-P2-PHY-S1-Q01, source={TOAS - Test 12 (10 Jul 2022)}, paper=2, qno=1, type=single, subtopic={Processes}, answer={A}, answer_src=key, difficulty=3, diff_parts={K1 S2 I1 C0 F0}, diff_note={concepts: internal energy along a linear P-V path | steps: equation of the line; U proportional to PV; maximum of PV; volume where PV is half | I: PV along a line is quadratic}]
One mole of an ideal gas is expanded from the state A $(P_0,V_0)$ to final state B having volume $V$. The process follows a path represented by a straight line on the P-V diagram as shown. Up to what volume $(V)$ the gas shall be expanded so that final internal energy is half the maximum internal energy during the process.
\bfig{76.5}{TOAS12_p2_phy_s1_q01_1.png}
\begin{qoptions}
\opt{A}{$\left(\dfrac{\sqrt2+1}{\sqrt2}\right)V_0$}
\opt{B}{$\sqrt2V_0$}
\opt{C}{$(\sqrt2+1)V_0$}
\opt{D}{$\left(\dfrac{\sqrt2+1}{2}\right)V_0$}
\end{qoptions}
\end{bankq}
\begin{bankpassage}[id=TOAS12-P2-PHY-S2-G11-12, source={TOAS - Test 12 (10 Jul 2022)}, questions={11-12}]
Two moles of an ideal mono-atomic gas is taken through a cyclic process ABCA as shown in the V-T diagram. In the process BC, $VT=$ constant. The temperature at A is $T_0=300$ K. [take $\ln(3)=1.1$]. Then answer the following questions.
\bfig{98.5}{TOAS12_p2_phy_s2_para11_12_1.png}
\end{bankpassage}
\begin{bankq}[id=TOAS12-P2-PHY-S2-Q11, source={TOAS - Test 12 (10 Jul 2022)}, paper=2, qno=11, type=multi, subtopic={Heat engines}, answer={C}, answer_src=key, difficulty=3, diff_parts={K1 S2 I0 C0 F1}, diff_note={concepts: work in a cyclic process | steps: state variables from the V-T diagram; process BC with VT constant; work by integration}, group=TOAS12-P2-PHY-S2-G11-12]
The work done by the gas during the process `BC' is :-
\begin{qoptions}
\opt{A}{600 R}
\opt{B}{900 R}
\opt{C}{1200 R}
\opt{D}{1500 R}
\end{qoptions}
\end{bankq}
\begin{bankq}[id=TOAS12-P2-PHY-S2-Q12, source={TOAS - Test 12 (10 Jul 2022)}, paper=2, qno=12, type=multi, subtopic={Heat engines}, answer={D}, answer_src=key, difficulty=4, diff_parts={K1 S2 I0 C1 F1}, diff_note={concepts: efficiency of a cyclic process | steps: heat absorbed and rejected in each leg; net work; efficiency}, group=TOAS12-P2-PHY-S2-G11-12]
The efficiency (in percentage) of the cycle is :-
\begin{qoptions}
\opt{A}{12}
\opt{B}{18}
\opt{C}{24}
\opt{D}{36}
\end{qoptions}
\end{bankq}
\begin{bankq}[id=TOAS13-P1-PHY-S1-Q05, source={TOAS - Test 13 (24 Jul 2022)}, paper=1, qno=5, type=multi, subtopic={First law}, answer={A,C}, answer_src=key, difficulty=2, diff_parts={K1 S1 I0 C0 F1}, diff_note={concepts: first law over a cycle | steps: delta U of the cycle is zero so the two changes cancel; Q = W for the cycle}]
In the cyclic process shown in the given figure, $\Delta U_1$ and $\Delta U_2$ represent the change in internal energy in process A and B, respectively. If $\Delta Q$ be the net heat given to the system in the process and $\Delta W$ be the work done by the system in the process, then
\bfig{62.0}{TOAS13_p1_phy_s1_q05_1.png}
\begin{qoptions}
\opt{A}{$\Delta U_1+\Delta U_2=0$}
\opt{B}{$\Delta U_1-\Delta U_2=0$}
\opt{C}{$\Delta Q-\Delta W=0$}
\opt{D}{$\Delta Q+\Delta W=0$}
\end{qoptions}
\end{bankq}
\begin{bankq}[id=TOAS13-P1-PHY-S1-Q06, source={TOAS - Test 13 (24 Jul 2022)}, paper=1, qno=6, type=multi, subtopic={Processes}, answer={A,B,C,D}, answer_src=key, difficulty=3, diff_parts={K1 S2 I0 C0 F1}, diff_note={concepts: internal energy of a monatomic gas | steps: read temperatures at a b c from PV = nRT; U = 3/2 nRT at each state; compare differences with the options}]
One mole of an ideal monatomic gas has initial temperature $T_0$, is made to go through the cycle abca shown in the given figure. If $U$ denotes the internal energy, then choose the correct alternative.
\bfig{111.1}{TOAS13_p1_phy_s1_q06_1.png}
\begin{qoptions}
\opt{A}{$U_c>U_b>U_a$}
\opt{B}{$U_c-U_b=3RT_0$}
\opt{C}{$U_c-U_a=\dfrac{9RT_0}{2}$}
\opt{D}{$U_b-U_a=\dfrac{3RT_0}{2}$}
\end{qoptions}
\end{bankq}
\begin{bankq}[id=TOAS13-P1-PHY-S3-Q07, source={TOAS - Test 13 (24 Jul 2022)}, paper=1, qno=7, type=numerical, subtopic={Processes}, answer={3.00}, answer_src=key, difficulty=4, diff_parts={K2 S2 I0 C0 F1}, diff_note={concepts: adiabatic law, isothermal law | steps: isothermal BA relates the pressures and volumes; adiabatic AC gives gamma; C\_V = R/(gamma - 1)}]
In the figure three process on a gas are shown. The value of $C_V$ is $\alpha R$ for the gas. If AC represents adiabatic process and BA represents isothermal process then $\alpha$ ?
\bfig{122.2}{TOAS13_p1_phy_s3_q07_1.png}
\end{bankq}
\begin{bankq}[id=TOAS13-P2-PHY-S1-Q07, source={TOAS - Test 13 (24 Jul 2022)}, paper=2, qno=7, type=single, subtopic={Processes}, answer={A}, answer_src=key, difficulty=1, diff_parts={K1 S1 I0 C0 F0}, diff_note={concepts: adiabatic law in terms of T and V | steps: 1/V proportional to T\^{}(1/(gamma-1)); parabola means exponent 2 so gamma = 3/2}]
Figure shows a parabolic curve between $\dfrac1V$ and $T$ for a mixture of a gas undergoing adiabatic process.
\bfig{113.7}{TOAS13_p2_phy_s1_q07_1.png}
The value of adiabatic exponent ($\gamma$) of mixture is
\begin{qoptions}
\opt{A}{$\dfrac32$}
\opt{B}{$\dfrac53$}
\opt{C}{$\dfrac75$}
\opt{D}{$\dfrac52$}
\end{qoptions}
\end{bankq}
\begin{bankq}[id=TOAS13-P2-PHY-S1-Q11, source={TOAS - Test 13 (24 Jul 2022)}, paper=2, qno=11, type=single, subtopic={First law}, answer={A}, answer_src=key, difficulty=1, diff_parts={K1 S1 I0 C0 F0}, diff_note={concepts: isobaric heat and work | steps: W/Q = R/Cp = 2/5 for a monatomic gas}]
100 J of heat is given to a monatomic gas sample slowly such that pressure remains constant. The work done by gas is :
\begin{qoptions}
\opt{A}{40 J}
\opt{B}{80 J}
\opt{C}{60 J}
\opt{D}{20 J}
\end{qoptions}
\end{bankq}
\begin{bankq}[id=TOAS13-P2-PHY-S1-Q17, source={TOAS - Test 13 (24 Jul 2022)}, paper=2, qno=17, type=single, subtopic={Processes}, answer={C}, answer_src=key, difficulty=2, diff_parts={K1 S2 I0 C0 F0}, diff_note={concepts: molar heat capacity of a polytropic process | steps: P proportional to V so x = -1; dQ = Cv dT + P dV; C = Cv + R/(1-x)}]
One mole of an ideal monatomic gas undergoes the process $P=2\sqrt T$. The molar specific heat of gas is
\begin{qoptions}
\opt{A}{$\dfrac{3R}{4}$}
\opt{B}{$\dfrac{9R}{4}$}
\opt{C}{2R}
\opt{D}{$\dfrac{2R}{3}$}
\end{qoptions}
\end{bankq}
\begin{bankq}[id=TOAS13-P2-PHY-S1-Q19, source={TOAS - Test 13 (24 Jul 2022)}, paper=2, qno=19, type=single, subtopic={Heat pumps}, answer={C}, answer_src=key, difficulty=1, diff_parts={K1 S1 I0 C0 F0}, diff_note={concepts: coefficient of performance of a heat pump | steps: COP = Th/(Th - Tc); Qh = COP times W}]
A heat pump operating between reservoirs at $-5^\circ$C and $20^\circ$C requires an electrical input of 1.2 kJ in each cycle. The heat dissipated into hot reservoir in each cycle.
\begin{qoptions}
\opt{A}{12.8 kJ}
\opt{B}{11.7 kJ}
\opt{C}{14 kJ}
\opt{D}{10 kJ}
\end{qoptions}
\end{bankq}
\begin{bankq}[id=TOAS13-P2-PHY-S2-Q09, source={TOAS - Test 13 (24 Jul 2022)}, paper=2, qno=9, type=numerical, subtopic={Entropy}, answer={489.00}, answer_src=key, difficulty=2, diff_parts={K1 S1 I0 C0 F1}, diff_note={concepts: entropy of a phase change | steps: delta S = m L/T}]
An ice cube of mass 400 g at temperature $0^\circ$C melts to water at $0^\circ$C. Heat is supplied from the surrounding air at slightly above $0^\circ$C. The process takes place very slowly. Find the change in entropy (J/K) of ice when it has all melted. (Latent heat of ice $=334$ kJ/kg)
\end{bankq}
\begin{bankq}[id=TOAS14-P1-PHY-S1-Q01, source={TOAS - Test 14 (07 Aug 2022)}, paper=1, qno=1, type=multi, subtopic={Entropy}, answer={B,C,D}, answer_src=key, difficulty=5, diff_parts={K2 S2 I0 C1 F1}, diff_note={concepts: work in a cyclic process, entropy of an ideal gas | steps: identify isobaric and isochoric legs from the graph; work as areas including the curved legs; entropy over a reversible cycle; delta S of the leg CD from the temperature and pressure ratios}]
2 mole diatomic gas undergoes cyclic process ABCDA as shown in indicator diagram. Then select correct alternative :-
\bfig{84.6}{TOAS14_p1_phy_s1_q01_1.png}
\begin{qoptions}
\opt{A}{Process AB, CD are isobaric and process BC, DA are isochoric.}
\opt{B}{Work done in a cycle is $100R\left\{5-\ln\left(\dfrac{16}{9}\right)\right\}$.}
\opt{C}{Change in entropy of gas and surrounding is zero in a cycle.}
\opt{D}{Change in entropy of gas in process CD is $R\ln\left\{\dfrac{(3)^2(4)^3}{(5)^5}\right\}$.}
\end{qoptions}
\end{bankq}
\begin{bankq}[id=TOAS19-P1-PHY-S1-Q03, source={TOAS - Test 19 (05 Feb 2023)}, paper=1, qno=3, type=multi, subtopic={Processes}, answer={A,B,C}, answer_src=key, difficulty=5, diff_parts={K1 S2 I1 C1 F1}, diff_note={concepts: polytropic process of an ideal gas | steps: heat capacity for P proportional to V; variable heat capacity for the general line; isobaric case; slope of temperature against volume}, also_in={TNAS - Test 01 (05 Feb 2023) P2 Q3}]
An ideal monoatomic gas is taken through a thermodynamics process $P=aV+b$ when $P$ and $V$ pressure and volume of gas and $a$ and $b$ are positive constants.
\begin{qoptions}
\opt{A}{If $b=0$, molar heat capacity of gas is constant.}
\opt{B}{If $a\neq0$, $b\neq0$, molar heat capacity of gas is lower in higher volume}
\opt{C}{If $a=0$, increase in volume per unit rise in temperature is constant.}
\opt{D}{Slope of temperature versus volume graph is constant.}
\end{qoptions}
\end{bankq}

\section*{TAAS papers}
\begin{bankq}[id=TAAS02-P1-PHY-S2-Q05, source={TAAS - Test 02 (30 Jul 2023)}, paper=1, qno=5, type=multi, subtopic={Processes}, answer={A,B}, answer_src=key, difficulty=3, diff_parts={K1 S2 I0 C0 F1}, diff_note={concepts: isothermal and adiabatic compression | steps: final temperature; final pressure; work for each; compare}]
Two gases have same initial pressure, volume, temperature. Both gases are compressed to same volume one isothermally and one adiabatically :-
\begin{qoptions}
\opt{A}{The final temperature is greater for the adiabatic process.}
\opt{B}{The final pressure is greater for the adiabatic process.}
\opt{C}{The work done on the gas is greater for the isothermal process.}
\opt{D}{All the above options are incorrect.}
\end{qoptions}
\end{bankq}
\begin{bankq}[id=TAAS02-P1-PHY-S2-Q06, source={TAAS - Test 02 (30 Jul 2023)}, paper=1, qno=6, type=multi, subtopic={Processes}, answer={A,D}, answer_src=key, difficulty=4, diff_parts={K1 S2 I1 C0 F1}, diff_note={concepts: cycle with isothermal and isochoric parts | steps: identify processes; translate to another diagram; match the options}]
A cyclic process $ABCD$ is shown in the $P$-$V$ diagram. Which of the following curves represents the same process if $BC$ and $DA$ are isothermal processes.
\bfig{71.1}{TAAS02_p1_phy_s2_q06_1.png}
\begin{qoptions}
\opt{A}{\bfigi{68.5}{TAAS02_p1_phy_s2_q06_A.png}}
\opt{B}{\bfigi{68.9}{TAAS02_p1_phy_s2_q06_B.png}}
\opt{C}{\bfigi{68.5}{TAAS02_p1_phy_s2_q06_C.png}}
\opt{D}{\bfigi{68.9}{TAAS02_p1_phy_s2_q06_D.png}}
\end{qoptions}
\end{bankq}
\begin{bankq}[id=TAAS02-P1-PHY-S4-Q01, source={TAAS - Test 02 (30 Jul 2023)}, paper=1, qno=1, type=numerical, subtopic={First law}, answer={3}, answer_src=key, difficulty=6, diff_parts={K2 S2 I1 C1 F1}, diff_note={concepts: heat in isothermal, isochoric processes | steps: work in the isothermal expansion; temperature after the isochoric heating; heat at constant volume; solve for n}]
$n$-moles of an ideal gas at $27^\circ$C are expanded isothermally to five time its volume and then heated at constant volume to raise its pressure equal to initial pressure (before expansion). If heat given to the system in the process is 83.14 kJ then find the value of $n$ (to the nearest integer). The $\gamma\left(=\dfrac{C_P}{C_V}\right)$ of the gas is 1.42. [$\ell$n5 $=1.61$]
\end{bankq}
\begin{bankq}[id=TAAS02-P1-PHY-S4-Q03, source={TAAS - Test 02 (30 Jul 2023)}, paper=1, qno=3, type=numerical, subtopic={First law}, answer={8}, answer_src=key, difficulty=6, diff_parts={K1 S2 I2 C1 F1}, diff_note={concepts: power balance with a moving piston | steps: electric power equals the rate of work at constant temperature; force from the weight; speed}]
A resistance coil wired to an external battery is placed inside a thermally insulated cylinder fitted with a frictionless piston and containing an ideal gas. A current $I=2.40$ mA flows through the coil, which has a resistance $R=490\ \Omega$. If the speed $v$ with which piston of mass $m=12$ kg move upward, so that the temperature of the gas remains unchanged is 0.003$x$ mm/sec. Find $x$?
\bfig{76.4}{TAAS02_p1_phy_s4_q03_1.png}
\end{bankq}

\section*{TOPS papers}
\begin{bankq}[id=TOPS05-P1-PHY-S1-Q04, source={TOPS - Test 05 (14 Jan 2023)}, paper=1, qno=4, type=multi, subtopic={Processes}, answer={B,C}, answer_src=key, difficulty=6, diff_parts={K2 S2 I1 C1 F1}, diff_note={concepts: isothermal curve, geometry | steps: slope of the curve; intercept of the tangent; coordinates of D; area of the triangle}]
$n$ moles of an ideal gas undergo an isothermal process at temperature $T$. $P$-$V$ graph of the process is as shown in the figure. A point $A(V_1,P_1)$ is located on the $P$-$V$ curve. Tangent at point $A$, cuts the $V$-axis at point $D$. $AO$ is the line joining the point $A$ to the origin $O$ of $PV$ diagram. Then,
\bfig{89.4}{TOPS05_p1_phy_s1_q04_1.png}
\begin{qoptions}
\opt{A}{coordinates of points $D$ is $\left(\dfrac{3V_1}2,0\right)$}
\opt{B}{coordinates of points $D$ is $(2V_1,0)$}
\opt{C}{area of the triangle $AOD$ is $nRT$}
\opt{D}{area of the triangle $AOD$ is $\dfrac34nRT$}
\end{qoptions}
\end{bankq}
\begin{bankq}[id=TOPS05-P1-PHY-S1-Q05, source={TOPS - Test 05 (14 Jan 2023)}, paper=1, qno=5, type=multi, subtopic={First law}, answer={A,B,C}, answer_src=key, difficulty=7, diff_parts={K3 S2 I1 C1 F1}, diff_note={concepts: polytropic process, work, heat capacity | steps: exponent of the process; work and change in internal energy; molar heat capacity}]
Two moles of a monatomic ideal gas undergoes a thermodynamic process $\dfrac{V^3}{T^2}=$ constant, if the temperature is raised by 300 K then
\begin{qoptions}
\opt{A}{work done by the gas is 400 $R$}
\opt{B}{change in internal energy is 900 $R$.}
\opt{C}{molar heat capacity of the gas for the process is 13/6 $R$}
\opt{D}{molar heat capacity of the gas for the process is 3/2 $R$.}
\end{qoptions}
\end{bankq}
\begin{bankq}[id=TOPS05-P1-PHY-S2-Q09, source={TOPS - Test 05 (14 Jan 2023)}, paper=1, qno=9, type=single, subtopic={Processes}, answer={A}, answer_src=key, difficulty=5, diff_parts={K2 S2 I1 C1 F0}, diff_note={concepts: thermodynamic processes, signs | steps: processes on the diagram; change in internal energy; heat and work signs; match}]
One mole of a monatomic gas is taken through a cycle $ABCDA$ as shown in the $P$-$V$ diagram. List--II give the characteristics involved in the cycle. Match them with each of the processes given in list--I.
\bfig{118.7}{TOPS05_p1_phy_s2_q09_1.png}
{\footnotesize\setlength{\tabcolsep}{3pt}
\begin{tabular}{|l|p{2.3cm}|l|p{2.3cm}|}
\hline
 & List--I & & List--II\\
\hline
(I) & Process $A\to B$ & (P) & Internal energy decreases\\
(II) & Process $B\to C$ & (Q) & Internal energy increases\\
(III) & Process $C\to D$ & (R) & Heat is lost\\
(IV) & Process $D\to A$ & (S) & Heat is gained\\
 & & (T) & Work is done on the gas\\
\hline
\end{tabular}}
\begin{qoptions}
\opt{A}{I$\to$P,R,T;II$\to$P,R;III$\to$Q,S;IV$\to$R,T}
\opt{B}{I$\to$P;II$\to$Q;III$\to$R,S,T;IV$\to$Q,T}
\opt{C}{I$\to$P,R,T;II$\to$P,R;III$\to$R,T;IV$\to$R,T}
\opt{D}{I$\to$R,T;II$\to$Q,S;III$\to$P;IV$\to$P,R,T}
\end{qoptions}
\end{bankq}
\begin{bankq}[id=TOPS05-P1-PHY-S3-Q06, source={TOPS - Test 05 (14 Jan 2023)}, paper=1, qno=6, type=numerical, subtopic={First law}, answer={900.00}, answer_src=key, difficulty=6, diff_parts={K2 S2 I1 C1 F1}, diff_note={concepts: adiabatic process, internal energy | steps: temperature relation; final temperature; decrease in energy}]
One mole of a monatomic ideal gas undergoes an adiabatic expansion in which its volume becomes eight times its initial value. If the initial temperature of the gas is 100 K and the universal gas constant $R=8.0$ J mol$^{-1}$ K$^{-1}$, the decrease in its internal energy. in Joule, is \underline{\hspace{1.5cm}}
\end{bankq}
\begin{bankq}[id=TOPS05-P2-PHY-S2-Q06, source={TOPS - Test 05 (14 Jan 2023)}, paper=2, qno=6, type=multi, subtopic={Processes}, answer={A,D}, answer_src=key, difficulty=5, diff_parts={K1 S2 I1 C1 F1}, diff_note={concepts: thermodynamic processes | steps: process slopes; internal energy; heat and isothermal conditions}]
An ideal gas is taken from thermodynamic state $A$ to $D$ as shown in the figure :-
\bfig{108.2}{TOPS05_p2_phy_s2_q06_1.png}
\begin{qoptions}
\opt{A}{Total change in internal energy is zero.}
\opt{B}{Process $CD$ is isobaric.}
\opt{C}{Process $BC$ may be isothermal.}
\opt{D}{In process $CD$ heat is released.}
\end{qoptions}
\end{bankq}

\section*{TEAS papers}
\begin{bankq}[id=TEAS01-P1-PHY-S3-Q07, source={TEAS - Test 01 (16 Apr 2023)}, paper=1, qno=7, type=numerical, subtopic={Entropy}, answer={2.00}, answer_src=key, difficulty=7, diff_parts={K2 S2 I2 C1 F1}, diff_note={concepts: entropy, heat capacity | steps: heat as T dS; differentiate; heat capacity expression; x}]
In a thermodynamic prcess, temperature of gas as function of entropy is given as $T=\alpha.S^2$ where `$T$' is temperature and `$S$' is entropy of gas `$\alpha$' is a positive constant. If heat capacity of gas as a function of entrpy is $C=\dfrac Sx$ then find `$x$'.
\end{bankq}

\section*{TNAS papers}
\begin{bankq}[id=TNAS01-P1-PHY-S1-Q08, source={TNAS - Test 01 (05 Feb 2023)}, paper=1, qno=8, type=single, subtopic={Processes}, answer={C}, answer_src=key, difficulty=4, diff_parts={K1 S2 I1 C1 F0}, diff_note={concepts: adiabatic process on log graph | steps: adiabatic relation; take logs; slope}]
Which among the following represents the slope of adiabatic process of an ideal gas in $\ell n T$ Vs $\ell n P$ graph.
\begin{qoptions}
\opt{A}{$\dfrac{C_P}{C_V}$}
\opt{B}{$1-\dfrac{C_P}{C_V}$}
\opt{C}{$1-\dfrac{C_V}{C_P}$}
\opt{D}{$-\dfrac{C_V}{C_P}$}
\end{qoptions}
\end{bankq}
\begin{bankq}[id=TNAS01-P1-PHY-S1-Q19, source={TNAS - Test 01 (05 Feb 2023)}, paper=1, qno=19, type=single, subtopic={Processes}, answer={C}, answer_src=key, difficulty=5, diff_parts={K1 S2 I2 C1 F0}, diff_note={concepts: process of a diatomic gas | steps: relation between P and V; work; change in internal energy; ratio}]
An ideal diatomic gas undergoes a process in which its internal energy $U$ and density $\rho$ vary as $U\rho=$ constant. The ratio of change in internal energy and the work done by the gas is :-
\begin{qoptions}
\opt{A}{$\dfrac32$}
\opt{B}{$\dfrac35$}
\opt{C}{$\dfrac52$}
\opt{D}{$\dfrac53$}
\end{qoptions}
\end{bankq}
\begin{bankq}[id=TNAS01-P1-PHY-S2-Q01, source={TNAS - Test 01 (05 Feb 2023)}, paper=1, qno=1, type=numerical, subtopic={Processes}, answer={12}, answer_src=key, difficulty=5, diff_parts={K1 S2 I1 C1 F1}, diff_note={concepts: molar heat capacity | steps: first law; work per degree; capacity}]
For an ideal gas $\left(\dfrac{C_P}{C_V}=\dfrac32\right)$ undergoing a thermodynamic process work done is related to temperature change as $dW=2RdT$ for one mole. Find molar heat capacity of the gas is $xR$, then value of $3x$ is :
\end{bankq}
\begin{bankq}[id=TNAS01-P1-PHY-S2-Q02, source={TNAS - Test 01 (05 Feb 2023)}, paper=1, qno=2, type=numerical, subtopic={Heat pumps}, answer={9}, answer_src=key, difficulty=3, diff_parts={K1 S1 I0 C1 F1}, diff_note={concepts: refrigerator coefficient | steps: coefficient of performance; heat from the cold reservoir}]
A Carnot heat engine, having efficiency $\eta=0.1$ is used as a refrigerator. If the work done on the system is $x$J, then the amount of energy absorbed from the cold reservoir is $N\,x$J, value of $N$ is
\end{bankq}
\begin{bankq}[id=TNAS01-P2-PHY-S3-Q04, source={TNAS - Test 01 (05 Feb 2023)}, paper=2, qno=4, type=numerical, subtopic={Heat engines}, answer={0.165}, answer_src=key, difficulty=6, diff_parts={K1 S2 I2 C1 F1}, diff_note={concepts: efficiency of a cycle | steps: processes; heat input; work; efficiency}]
Find the efficiency (in percentage) of the cycle shown in the diagram $p$, $V$ on figure. The working fluid is one mole of an ideal monatomic gas.
\bfig{100.1}{TNAS01_p2_phy_s3_q04_1.png}
\end{bankq}
\begin{bankq}[id=TNAS01-P2-PHY-S3-Q05, source={TNAS - Test 01 (05 Feb 2023)}, paper=2, qno=5, type=numerical, subtopic={First law}, answer={1133.33}, answer_src=key, difficulty=5, diff_parts={K1 S2 I1 C1 F1}, diff_note={concepts: work from a cycle table | steps: energy balance of each process; net work; rate}]
A fluid system, contained in a piston and cylinder machine, passes through a complete cycle of four processes. The sum of all heat transferred during a cycle is $-340$ kJ. The system completes $200$ cycles per min. The following table shows the heat, work and internal energy for each process. The rate of work output (in kW) is
\begin{center}
{\footnotesize\setlength{\tabcolsep}{2pt}\begin{tabular}{|l|l|l|l|}\hline
Process & $Q$ (kJ/min) & $W$ (kJ/min) & $\Delta E$ (kJ/min)\\\hline
1-2 & $0$ & $4340$ & $-$\\\hline
2-3 & $42000$ & $0$ & $-$\\\hline
3-4 & $-4200$ & $-$ & $-73200$\\\hline
4-1 & $-$ & $-$ & $-$\\\hline
\end{tabular}}
\end{center}
\end{bankq}
\begin{bankq}[id=TNAS01-P2-PHY-S3-Q06, source={TNAS - Test 01 (05 Feb 2023)}, paper=2, qno=6, type=numerical, subtopic={Entropy}, answer={4.275}, answer_src=key, difficulty=6, diff_parts={K1 S2 I2 C1 F1}, diff_note={concepts: entropy change with variable specific heat | steps: integrate cv over T; mass; value}]
Determine the entropy change of $4$ kg of a perfect gas whose temperature varies from $127^\circ$C to $227^\circ$C during a constant volume process. The specific heat varies linearly with absolute temperature and is represented by the relation : $c_v=(0.48+0.0096\,T)$ kJ/kg K. [Take $\ell$n $(1.25)=0.22$]
\end{bankq}
\begin{bankq}[id=TNAS05-P1-PHY-S1-Q08, source={TNAS - Test 05 (09 Jul 2023)}, paper=1, qno=8, type=single, subtopic={Entropy}, answer={B}, answer_src=key, difficulty=6, diff_parts={K2 S2 I1 C2 F0}, diff_note={concepts: polytropic process, entropy change of a gas | steps: molar heat capacity of the process; temperature ratio; entropy integral; evaluate}]
One mole of an ideal gas with adiabatic exponent $\gamma=1.5$ goes through a polytropic process at the end of which the absolute temperature of gas increases 2 times. The polytropic constant of gas equal 2 then increase in entropy of gas in this process. (In SI units)
\begin{qoptions}
\opt{A}{4.76}
\opt{B}{5.76}
\opt{C}{4.26}
\opt{D}{5.26}
\end{qoptions}
\end{bankq}
\begin{bankq}[id=TNAS05-P2-PHY-S1-Q03, source={TNAS - Test 05 (09 Jul 2023)}, paper=2, qno=3, type=multi, subtopic={Processes}, answer={A,B,D}, answer_src=key, difficulty=8, diff_parts={K1 S3 I2 C2 F1}, diff_note={concepts: ideal gas on two sides of a piston with springs | steps: pressure on the right; isothermal left side; spring energy; work on the piston; temperature; internal energy; heat}]
A horizontal tank of length $2L$ is divided into two equal parts by a thermally insulating piston. Both parts contain a sample of ideal diatomic gas of $n$ moles each at temperature $T_0$. To each side of piston a ideal spring of spring constant $k$ is attached such that the other ends of spring are connected by the walls of the tank. The springs are in natural length when the piston is at the center of the tank. When the gas in right part is slowly heated, the piston moves a distance $\dfrac L2$ towards the left. During the process, the left part is connected to a heat reservoir at temperature $T_0$. Given $k=\dfrac{nRT_0}{L^2}$.
\begin{qoptions}
\opt{A}{The pressure of the gas in the right part when the piston has moved a distance $\dfrac L4$ from the center is $\dfrac{11}{6}\dfrac{nRT_0}{AL}$.}
\opt{B}{The pressure of the gas in the right part when the piston has moved a distance $\dfrac L2$ from the center is $\dfrac{3nRT_0}{AL}$.}
\opt{C}{The work done by gas in the right part on the piston is $0.85\ nRT_0$.}
\opt{D}{The total heat supplied to the right part of the gas is $9.7\ nRT_0$.}
\end{qoptions}
\end{bankq}
\begin{bankq}[id=TNAS05-P2-PHY-S1-Q06, source={TNAS - Test 05 (09 Jul 2023)}, paper=2, qno=6, type=multi, subtopic={Heat engines}, answer={A,D}, answer_src=key, difficulty=8, diff_parts={K1 S3 I2 C2 F1}, diff_note={concepts: Carnot engine and heat pump with objects of finite heat capacity | steps: engine between two objects; final temperature; work; refrigerator between the others; energy balance; compare}]
Consider three objects with equal and constant heat capacity at initial temperatures $T_1=225$ K, $T_2=400$ K and $T_3=375$ K. Suppose we wish to raise the temperature of the third object. To do this, we could run a Carnot engine between the first and second objects, extracting work $W$. This work can then be dissipated as heat to raise the temperature of the third object. Even better. it can be stored and used to run a Carnot engine between the first and third object in reverse, which pumps heat into the third object. Assume that all work produced by running engines can be stored and used without dissipation.
\begin{qoptions}
\opt{A}{The minimum temperature $T_L$ to which the first object can be lowered is 300 K.}
\opt{B}{The minimum  temperature $T_2$ to which the first object can be lowered is 250 K.}
\opt{C}{The maximum temperature $T_H$ to which the third object can be raised is 400 K.}
\opt{D}{The maximum temperature $T_H$ to which the third object can be raised is 450 K.}
\end{qoptions}
\end{bankq}

\section*{Nurture (General) papers}
\begin{bankq}[id=NUR01-P2-PHY-S3-Q03, source={Nurture (General) - Test 01 (12 Jun 2021)}, paper=2, qno=3, type=single, subtopic={Heat engines}, answer={A}, answer_src=solved, working={Otto cycle efficiency 1-(1/r)\^{}(gamma-1)}, difficulty=2, diff_parts={K1 S1 I1 C0 F0}, diff_note={concepts: efficiency of the Otto cycle | steps: work and heat; temperature ratios}]
For a given mass of working gas of adiabatic constant $\gamma\left(=\dfrac{C_P}{C_V}\right)$ if the compression ratio $r$ is defined as $r=\dfrac{V_1}{V_2}$ (Where $V_1$ is volume of gas at 1 and $V_2$ is volume of gas at 2). Then the efficiency of the cycle is given as.
\begin{qoptions}
\opt{A}{$\eta_{th}=1-\left(\dfrac1r\right)^{\gamma-1}$}
\opt{B}{$\eta_{th}=1+\left(\dfrac1r\right)^{\gamma-1}$}
\opt{C}{$\eta_{th}=1+\left(\dfrac1r\right)^{\gamma+1}$}
\opt{D}{$\eta_{th}=1-\left(\dfrac1r\right)^{\gamma+1}$}
\end{qoptions}
\end{bankq}
\begin{bankq}[id=NUR01-P2-PHY-S3-Q04, source={Nurture (General) - Test 01 (12 Jun 2021)}, paper=2, qno=4, type=single, subtopic={Heat engines}, answer={A}, answer_src=solved, working={Net work divided by the displacement volume V1(r-1)/r}, difficulty=6, diff_parts={K2 S2 I2 C1 F0}, diff_note={concepts: mean effective pressure of the Otto cycle, heat engines | steps: net work; number of moles; displacement volume; expression}]
If the displacement volume is given by $(V_1-V_2)$, then the mean effective pressure of the above cycle is given as
\begin{qoptions}
\opt{A}{$\left(\dfrac1{\gamma-1}\right)\left(\dfrac{P_1}{T_1}\right)\left(\dfrac r{r-1}\right)\{(T_3-T_2)-(T_4-T_1)\}$}
\opt{B}{$\left(\dfrac1{\gamma+1}\right)\left(\dfrac{P_1}{T_1}\right)\left(\dfrac r{r-1}\right)\{(T_3-T_2)+(T_4-T_1)\}$}
\opt{C}{$\left(\dfrac1{\gamma-1}\right)\left(\dfrac{P_1}{T_1}\right)\left(\dfrac r{r+1}\right)\{(T_3-T_2)+(T_4-T_1)\}$}
\opt{D}{$\left(\dfrac1{\gamma-1}\right)\left(\dfrac{P_1}{T_1}\right)\left(\dfrac r{r-1}\right)\{(T_3+T_2)+(T_4-T_1)\}$}
\end{qoptions}
\end{bankq}
\begin{bankq}[id=NUR04-P1-PHY-S3-Q18, source={Nurture (General) - Test 04 (11 Feb 2024)}, paper=1, qno=18, type=numerical, subtopic={Processes}, answer={200}, answer_src=solved, working={V\^{}2T constant gives PV\^{}3 constant; W=R(T1-T2)/(3-1)=-200R}, difficulty=7, diff_parts={K2 S2 I2 C1 F1}, diff_note={concepts: work in a polytropic process, processes | steps: convert to the polytropic exponent; work formula; evaluate}]
1 mole of an ideal monoatomic gas is expanded till the temperature of the gas is doubled under the process $V^2T$ = constant. The initial temperature of the gas is 400 K. In terms of R, the total work done in the process is -nR find n
\end{bankq}
\end{document}
